\epanalipsi{Γ1}

\tmima{Λεξικό του κεφαλαίου}

\renewcommand{\arraystretch}{1.5}
\begin{longtable}{@{}>{\raggedright\arraybackslash}p{0.28\linewidth} >{\raggedright\arraybackslash}p{0.62\linewidth}@{}}
\toprule
\textbf{Λέξη} & \textbf{Με μία πρόταση}\\
\midrule
\endhead
πληθυσμός, μεταβλητή & Αυτό που μελετάμε και το χαρακτηριστικό του που μας ενδιαφέρει.\\
απογραφή, δείγμα & Εξετάζω όλο τον πληθυσμό· εξετάζω ένα μέρος του.\\
αντιπροσωπευτικό δείγμα & Μικρογραφία του πληθυσμού, που δίνει αξιόπιστα συμπεράσματα.\\
κατηγορικά, ποσοτικά δεδομένα & Δεν είναι αριθμοί· είναι αριθμοί.\\
διακριτά, συνεχή & Μεμονωμένες τιμές (πόσα)· οποιαδήποτε τιμή σε διάστημα (πόσο).\\
συχνότητα, σχετική συχνότητα & Πόσες φορές εμφανίζεται μια τιμή· συχνότητα διά το σύνολο, ως ποσοστό.\\
κλάση, πλάτος, κέντρο & Ομάδα τιμών· διαφορά των άκρων της· μέσο των άκρων της.\\
ραβδόγραμμα & Ράβδοι με κενά, ύψος ίσο με τη συχνότητα. Για κατηγορικά και διακριτά.\\
κυκλικό διάγραμμα & Τομείς με $\hat{\omega}=\frac{\text{συχνότητα}}{\text{σύνολο}}\cdot360^\circ$.\\
ιστόγραμμα & Κολλητά ορθογώνια με βάση το πλάτος της κλάσης και ύψος τη συχνότητά της.\\
\bottomrule
\end{longtable}

\tmima{Ασκήσεις επανάληψης}

\ex Κατηγορικά ή ποσοτικά δεδομένα; Στα ποσοτικά πες αν είναι διακριτά ή συνεχή. α) Ο αριθμός των κατοικίδιων μιας οικογένειας. β) Το χρώμα της μπλούζας. γ) Ο χρόνος κολύμβησης στα $100$ μέτρα. δ) Το ύψος ενός δέντρου. ε) Το πλήθος των μαθητών ενός τμήματος. στ) Η αγαπημένη εφαρμογή στο κινητό.

\ex Η Αγγελική θέλει να μάθει τι είδους μουσική ακούν οι $300$ μαθητές του σχολείου της και ρωτά $40$ παιδιά που επιλέγει με κλήρωση από όλες τις τάξεις. Ποιος είναι ο πληθυσμός, ποιο το δείγμα, ποια η μεταβλητή και τι είδους είναι τα δεδομένα;

\ex Ο Νίκος θέλει να μάθει πόσα χρήματα ξοδεύουν οι μαθητές του σχολείου στο κυλικείο και ρωτά τα παιδιά που περιμένουν στην ουρά του κυλικείου. Είναι το δείγμα του αντιπροσωπευτικό; Δικαιολόγησε.

\ex Η Ζωή ρώτησε $20$ συμμαθητές ποιο είναι το αγαπημένο τους φρούτο: Μήλο (Μ), Βερίκοκο (Β), Πορτοκάλι (Π), Φράουλα (Φ), Κεράσι (Κ).\par
\noindent Β Π Κ Μ Π Π Β Φ Φ Μ Μ Μ Β Β Μ Κ Β Φ Π Μ\par
Φτιάξε πίνακα συχνοτήτων με σχετικές συχνότητες και ραβδόγραμμα.

\ex Το μήκος άλματος (σε cm) $20$ παιδιών στο μάθημα της γυμναστικής:\par
\noindent\mbox{}\hfill 151 \ 160 \ 138 \ 150 \ 172 \ 158 \ 142 \ 176 \ 167 \ 140 \ 170 \ 161 \ 181 \ 147 \ 188 \ 165 \ 145 \ 132 \ 155 \ 153\hfill\mbox{}\par
Ομαδοποίησε τα δεδομένα σε κλάσεις πλάτους $10$ ξεκινώντας από το $130$ και φτιάξε πίνακα συχνοτήτων και σχετικών συχνοτήτων.

\ex Για τις κλάσεις της προηγούμενης άσκησης βρες το πλάτος και το κέντρο καθεμιάς.

\ex Συμπλήρωσε τον πίνακα για $40$ παρατηρήσεις.
\begin{center}
\renewcommand{\arraystretch}{1.2}
\begin{tabular}{ccc}
\toprule
\textbf{Κλάσεις} & \textbf{Συχνότητα} & \textbf{Σχετική συχνότητα}\\
\midrule
$0-10$ & $4$ & $\ldots$\\
$10-20$ & $\ldots$ & $30\%$\\
$20-30$ & $14$ & $\ldots$\\
$30-40$ & $\ldots$ & $20\%$\\
$40-50$ & $2$ & $\ldots$\\
\midrule
\textbf{Σύνολο} & $40$ & $\ldots$\\
\bottomrule
\end{tabular}
\end{center}

\ex Με τα δεδομένα του πίνακα της άσκησης $\text{Ε}\Gamma1.5$ βρες πόσα παιδιά πήδηξαν τουλάχιστον $160$ cm και ποιο ποσοστό πήδηξε λιγότερο από $150$ cm.

\ex Ρωτήσαμε $24$ παιδιά ποιο παιχνίδι προτιμούν στο διάλειμμα: κρυφτό $8$, μπάλα $6$, σχοινάκι $6$, κυνηγητό $4$. Βρες τις επίκεντρες γωνίες και φτιάξε το κυκλικό διάγραμμα.

\ex Σε κυκλικό διάγραμμα για $72$ μαθητές οι γωνίες των τομέων είναι: βόλεϊ $100^\circ$, μπάσκετ $80^\circ$, τένις $60^\circ$, ποδόσφαιρο $120^\circ$. Πόσοι μαθητές προτιμούν κάθε άθλημα;

\ex Οι $25$ βαθμοί της κυρίας Αλεξάνδρας στην παράγραφο Γ1.2 έδωσαν συχνότητες $2$, $8$, $10$, $5$ στις κλάσεις $4-8$, $8-12$, $12-16$, $16-20$. Φτιάξε το ιστόγραμμα συχνοτήτων και το ιστόγραμμα σχετικών συχνοτήτων.

\ex Στο ιστόγραμμα φαίνονται τα λεπτά μελέτης την ημέρα $50$ παιδιών.
\begin{center}
\gaIstogramma{0}{20}{20}{5}{5,12,18,10,5}{μελέτη (λεπτά)}{παιδιά}
\end{center}
α) Πόσα παιδιά μελετούν λιγότερο από $40$ λεπτά και ποιο ποσοστό είναι αυτό; β) Πόσα μελετούν τουλάχιστον $60$ λεπτά; γ) Ποια κλάση έχει τη μεγαλύτερη συχνότητα;

\ex Διάλεξε το κατάλληλο διάγραμμα (ραβδόγραμμα, κυκλικό ή ιστόγραμμα) και εξήγησε: α) τα βάρη $150$ μαθητών, ομαδοποιημένα σε κλάσεις, β) τα μερίδια τεσσάρων εφαρμογών σε ένα σύνολο $100\%$, γ) πόσα παιδιά προτιμούν καθένα από $7$ είδη παγωτού.

\ex Η Πολυξένη βρήκε για τέσσερις κατηγορίες σχετικές συχνότητες $20\%$, $30\%$, $25\%$ και $15\%$ και έφτιαξε κυκλικό διάγραμμα. Τι δεν πάει καλά; Αν το λάθος είναι στην τρίτη κατηγορία, ποιες είναι οι σωστές γωνίες των τομέων;

\ex Σωστό ή λάθος; α) Τα συνεχή δεδομένα παίρνουν μεμονωμένες τιμές. β) Ο αριθμός των αδελφών ενός παιδιού είναι ποσοτικό διακριτό δεδομένο. γ) Η απογραφή εξετάζει όλα τα στοιχεία του πληθυσμού. δ) Ένα αντιπροσωπευτικό δείγμα δίνει αξιόπιστα συμπεράσματα. ε) Το χρώμα ματιών είναι ποσοτικό. στ) Η συχνότητα μιας τιμής είναι πάντα μικρότερη από το σύνολο ή ίση με αυτό. ζ) Το άθροισμα των σχετικών συχνοτήτων είναι πάντα $100\%$. η) Το δεξί άκρο μιας κλάσης ανήκει στην ίδια κλάση.

\ex Σωστό ή λάθος; α) Στο κυκλικό διάγραμμα ένα ποσοστό $25\%$ αντιστοιχεί σε γωνία $90^\circ$. β) Στο ιστόγραμμα τα ορθογώνια έχουν κενά. γ) Σε ένα ιστόγραμμα με κλάσεις ίσου πλάτους, το ύψος κάθε ορθογωνίου είναι η συχνότητα της κλάσης. δ) Το κέντρο της κλάσης $10-20$ είναι $15$. ε) Το κυκλικό διάγραμμα ταιριάζει σε πολλές κατηγορίες. στ) Στο ραβδόγραμμα οι ράβδοι ακουμπούν. ζ) Το ιστόγραμμα σχετικών συχνοτήτων έχει το ίδιο σχήμα με το ιστόγραμμα συχνοτήτων των ίδιων δεδομένων. η) Οι επίκεντρες γωνίες ενός κυκλικού διαγράμματος έχουν άθροισμα $180^\circ$.

\Needspace{17\baselineskip}\tmima{Ο χάρτης της επανάληψης}

Μετά τις λύσεις, δες σε ποιες ασκήσεις δυσκολεύτηκες. Ο πίνακας σου λέει ποια παράγραφο να ξαναδιαβάσεις.

\renewcommand{\arraystretch}{1.5}
\begin{tabular*}{\linewidth}{@{\extracolsep{\fill}} >{\raggedright\arraybackslash}p{0.38\linewidth} >{\raggedright\arraybackslash}p{0.52\linewidth}}
\toprule
\textbf{Αν δυσκολεύτηκες στις} & \textbf{ξαναδιάβασε τις παραγράφους}\\
\midrule
ΕΓ1.1, ΕΓ1.2, ΕΓ1.3, ΕΓ1.4, ΕΓ1.15 & 1.1 (πληθυσμός, δεδομένα, ραβδόγραμμα)\\
ΕΓ1.5, ΕΓ1.6, ΕΓ1.7, ΕΓ1.8 & 1.2 (κλάσεις και πίνακες)\\
ΕΓ1.9, ΕΓ1.10, ΕΓ1.11, ΕΓ1.12, ΕΓ1.13, ΕΓ1.14, ΕΓ1.16 & 1.3 (κυκλικά διαγράμματα και ιστογράμματα)\\
\bottomrule
\end{tabular*}
